\section{Discussion}
\label{sec:discussion}

\paragraph{When does set selection help?} Our results point to three conditions. The index must expose the evidence: on MemEye, almost every question word that occurs in an evidence round is visible in its index line, because images are the subject of their rounds and messages are short, whereas on MemLens most of these words are hidden in long messages about other topics, and the selector opens too few rounds. The question must refer to a set whose members a similarity score cannot rank together, which is where model judgment with a recall-oriented instruction adds recall over similarity and where a set output sizes the budget per question. Finally, the reader must profit from a small input: the lightweight reader gains from focused evidence on MemEye, whereas the strong reader on MemLens does better with 40 retrieved rounds than with the annotated evidence itself. When these conditions do not hold, wide similarity retrieval over the full round text is the safer default for a strong reader, and full-context reading remains its most accurate option on histories that fit in context. For a lightweight reader, the mechanism study points to a different remedy: rather than a more complete set, it needs a precise set whose members are made explicit, which is what the notes of VGS provide. The external test bounds this remedy as well: WorldMemArena hides as much of its evidence from the index as MemLens, yet the notes do not help, possibly because the reader there already performs well with wide retrieval and the notes summarize many members instead of a few. Across the four benchmarks, the gain of the set path follows the weakness of top-40 retrieval for the lightweight reader rather than the modality of individual questions: it is large where top-40 retrieval leaves this reader below 0.5 (MemEye and MemLens) and small where it exceeds 0.75 (Mem-Gallery and WorldMemArena), and within MemEye and MemLens it is as large on the questions that a question-only classifier marks as text-only ($+0.071$ and $+0.075$) as on the visual ones ($+0.039$ and $+0.062$). Index visibility therefore predicts when the selector can find the evidence, as the masking experiment shows, but whether the reader also needs the set made explicit depends on how the reader fails, which visibility does not measure.

\paragraph{Query time versus write time.} Across our experiments, write-time transformations of memory content lost information whenever they replaced the original, as dense or short captions instead of images and predicted groups did. Write-time artifacts helped when they served as navigation keys that leave the content intact, as captions in retrieval keys, index lines, and the compact index do. The held-out test adds that the fidelity of these keys bounds query-time selection, since a selector cannot resolve a reference that its index does not show. Deciding at question time which part of a long round the index should expose is therefore a natural next step, which we leave untested because designing it on MemLens would remove its status as held-out data.

\paragraph{Where to spend model capability.} The selection call reads text only, and its input changes only by appending, so strong-model capability spent on navigation is inexpensive when the index is cached, as in a persistent per-user memory: with an explicit context cache, the measured selection call on MemEye costs a fifth of the uncached call, whereas implicit caching hardly applied, and the compact index halves the input that remains. The reuse must be large enough: the cached strong pipeline saves money per call but pays to create and store each cache, so on WorldMemArena, where a cached index serves 11 questions, it costs more than top-40 retrieval, and roughly 17 questions per cached index are needed to break even. Without reuse, the strong selection call dominates the cost, and neither cheaper selectors nor smaller candidate pools preserved its quality. In either case, two properties of the strong selector remain: its output and thinking, which set a floor of about \$0.7 per 1,000 questions, and its decoding variance, which changes the selected sets between draws. Selection quality, moreover, did not follow model size or reasoning ability: \texttt{gpt-5} and a reasoning version of Qwen3 returned fewer rounds than \texttt{gemini-3.8-flash} under the same prompt. Selector prompts thus need calibration toward recall for each model family, and evaluations of index-based memory should report the number of opened rounds together with full-set recall.

\paragraph{Implications for evaluation.} Retrieval for memory should be reported with full-set recall, since recall@$k$ and coverage remain high while multi-item questions fail. Our results suggest that benchmark reports would benefit from two further measurements, index visibility (or, more generally, how much of the evidence a compact representation can expose) and the accuracy that survives the removal of all visual information. These properties separate MemEye (visibility 0.95; 0.191 without visual information) from MemLens (visibility 0.37) and Mem-Gallery (0.892 without visual information), and they determine which memory designs can help. Finally, benchmarks built from a few long histories need scenario-level uncertainty estimates, which in our case turned a significant question-level gain into an interval that includes zero.

\section{Limitations}
\label{sec:limitations}

The first group of limitations concerns the data. The development benchmark covers eight scenario histories whose questions follow recurring patterns of reference, and all selection mechanisms, prompts, and lightweight variants were designed on it; the lightweight variants are exploratory, and the first held-out test covers only the main conditions. MemLens differs from MemEye in several properties at once, such as message length, evidence style, and annotation granularity. Masking the question words in the MemEye index shows that lower visibility causes lower recall and accuracy, but this manipulation removes words rather than whole asides and leaves semantic cues, so it does not determine how much of the contrast between the two benchmarks visibility explains. Our gold rounds for MemLens are derived from session-level answer annotations and turn-level evidence flags, and the strong reader exceeds the annotated evidence, so full-set recall is an imperfect proxy on this benchmark. The scale test adds distractors from other users, which are easier to reject than distractors from the same user, so it does not establish robustness to confusable items of one history.

The second group concerns models and measurement. The readers, the main selector, the router, the note writer, and the captioner are Gemini models. Selectors from other families were tested only with the prompt that we developed with Gemini, so our results characterize this prompt-model combination rather than set selection in general, and a prompt calibrated toward recall for each family may change the ranking of selectors. The registration named Qwen3-235B-A22B as the cross-family selector on MemLens; we report it in the selection and end-to-end results and added \texttt{gpt-5-mini} after the registered runs. The judges, although from two families with high agreement, are language models rather than human annotators. The strong selector decodes nondeterministically, and four selection draws changed the accuracy of the lightweight reader without thinking by up to 0.034, which is comparable to some of the differences that we report for this reader.

The third group concerns the protocol and the cost estimates. The positive test of VGS uses a held-out part of the same benchmark, the aggregate MemLens results of the first test motivated the mechanism, and the gate threshold was set from two development distributions; the external test on WorldMemArena did not confirm the gain, and it covers only the lifelong-personal part of that benchmark, evaluated end to end because its evidence is annotated per session. The confirmation of RQ6 reuses the MemLens test-split questions with new, twice longer haystacks, and the abstention fallback was motivated by observations on the 32k test. Costs use the list prices and caching terms of one provider, and the cached costs were measured on two persistent benchmarks; other pricing schemes would change the absolute numbers, although the strong selection call dominates the cost under the prices that we used. Finally, we did not train selectors or readers, so we cannot say how much of the capability gap distillation would close.

\section{Ethical Considerations}
\label{sec:ethics}

The setting that we study, a persistent memory of a personal assistant, stores raw conversations, images, dates, and session structure of one user, and our mechanisms add derived artifacts: image descriptions, index lines, resolved questions with their evidence sets, and, for the cache-native index, a copy of the index held by the model provider for the lifetime of a context cache. Each of these artifacts can expose personal information more broadly than the original exchange, for example when a description names a person or a document in a photo, and each must follow the retention and deletion rules of the raw history: deleting a round has to remove its index line, any resolution that refers to it, and the cached index, which must then be rebuilt. Explicit caches expire after a set time and are scoped to one account, but per-user isolation, access control, and the consent of the user to provider-side storage remain obligations of the deploying system rather than properties of the method. A second risk is silent failure: an incomplete evidence set can yield a confident but wrong count or comparison, and neither \qss{} nor VGS detects missing members; the abstention fallback of Section~\ref{sec:rq6} covers only answers that the reader itself marks as unsupported. Finally, the strongest configurations rely on commercial models and caching features, which limits access and long-term reproducibility. All experiments in this paper use public research benchmarks and involve no new data from human participants.

\section{Conclusion}
\label{sec:conclusion}

We asked what remains unsolved when simple vector retrieval is already strong. On text benchmarks, raw retrieval with a large budget is statistically indistinguishable from reading the full history. In multimodal agent memory, questions often refer to sets of past items, and top-$k$ similarity rarely returns such sets completely. Query-time set selection over a lightweight index recovers complete sets with small, adaptive budgets on the benchmark where we developed it; same-model controls trace this recall mainly to model judgment and a recall-oriented instruction, with a smaller contribution from returning a set of model-chosen size. A pre-registered held-out test shows that these retrieval gains do not translate into better answers when the index hides the evidence and the reader is strong enough to use wide retrieval. Acting on this condition at question time, visibility-gated set reading widens the index and gives lightweight readers explicit per-member notes, which improves their answers on held-out MemLens questions where plain set selection did not, but not on an external benchmark with equally low visibility. A compact index cached once per user makes set selection with a strong reader both more accurate and cheaper than top-40 retrieval on the development benchmark when each cached index serves many questions, whereas single-use histories leave the cost of selection to be paid per question. Query-time set selection is therefore a conditional tool whose applicability can be estimated and acted upon at question time, rather than a replacement for similarity retrieval.
